\chapter{Graphs of the Trigonometric Functions}
\label{ch:ch06-graphs}

\section{From Motion to Graphs}

The previous chapter showed that sine and cosine arise naturally from
circular motion: a point moving around the unit circle with angular
velocity $\omega$ has coordinates $x(t) = \cos(\omega t)$ and $y(t) =
\sin(\omega t)$, equations that describe a position on a circle rather than
a curve on a page. A natural question follows immediately. What happens if
the circle itself is set aside and only one coordinate --- say, the
vertical one --- is tracked through time? As the point rises and falls
around the circle, its height changes continuously, and plotting angle
along the horizontal axis against height along the vertical axis produces
the familiar wave-shaped curve known as the graph of the sine function.
That curve is therefore not an arbitrary construction handed down by
convention; it is the geometric record of circular motion, unrolled onto a
flat page.

\section{Building the Sine Curve}

Consider a point moving around the unit circle, and recall the special
angles of Chapter~4. At $\theta = 0$ the point is $(1,0)$, so $\sin(0) = 0$;
at $\theta = \pi/2$ the point is $(0,1)$, so $\sin(\pi/2) = 1$; at $\theta =
\pi$ the point is $(-1,0)$, so $\sin(\pi) = 0$; at $\theta = 3\pi/2$ the
point is $(0,-1)$, so $\sin(3\pi/2) = -1$; and at $\theta = 2\pi$ the point
returns to $(1,0)$, so $\sin(2\pi) = 0$. Plotting these five points and
connecting them smoothly produces the sine curve, a graph that records
nothing more than the changing vertical position of the point rotating
around the circle in the previous chapter.

\section{Building the Cosine Curve}

The cosine function records the horizontal coordinate of the same rotating
point. At the same five special angles, $\cos(0) = 1$, $\cos(\pi/2) = 0$,
$\cos(\pi) = -1$, $\cos(3\pi/2) = 0$, and $\cos(2\pi) = 1$. The resulting
graph has the same wave shape as the sine curve but begins at its maximum
value rather than at zero, and like the sine graph it is a direct
consequence of the same circular motion, differing only in which coordinate
is being tracked.

\section{Periodic Functions}

Both graphs reveal a striking property: after one complete revolution, the
same pattern repeats exactly, a fact that is nothing more than the
coterminal-angle theorem of Chapter~4 seen from a new angle.

\begin{definition}
A function $f$ is \emph{periodic} if there exists a positive number $P$
such that $f(x+P) = f(x)$ for every $x$ in the domain of $f$. The smallest
such positive number is called the \emph{period} of $f$.
\end{definition}

\begin{definition}
The \emph{amplitude} of a periodic function is the maximum distance from
its midline to a peak or a trough. Amplitude measures the size of an
oscillation, while period measures the length of one complete cycle.
\end{definition}

\section{Period and Amplitude of Sine and Cosine}

The unit circle immediately supplies both quantities for sine and cosine.

\begin{theorem}
The functions $y = \sin\theta$ and $y = \cos\theta$ have period $2\pi$ and
amplitude $1$.
\end{theorem}

\begin{proof}
By the coterminal-angle theorem of Chapter~4, $\theta$ and $\theta + 2\pi$
determine the same point on the unit circle, so $\sin(\theta+2\pi) =
\sin\theta$ and $\cos(\theta+2\pi) = \cos\theta$; the period is therefore
$2\pi$. Every point $(x,y)$ on the unit circle satisfies $-1 \le x \le 1$
and $-1 \le y \le 1$, so the largest value either coordinate can attain is
$1$ and the smallest is $-1$. The distance from the midline $y=0$ to either
extreme is therefore $1$, which is the amplitude.
\end{proof}

\section{Symmetry}

The graphs possess symmetries that follow directly from the geometry of the
unit circle rather than from inspection of the curve alone.

\begin{theorem}
The sine function is odd: $\sin(-\theta) = -\sin\theta$.
\end{theorem}

\begin{proof}
Reflecting an angle across the $x$-axis reverses the sign of its vertical
coordinate while leaving the horizontal coordinate unchanged. Since sine
measures the vertical coordinate of a point on the unit circle,
$\sin(-\theta) = -\sin\theta$.
\end{proof}

\begin{theorem}
The cosine function is even: $\cos(-\theta) = \cos\theta$.
\end{theorem}

\begin{proof}
Reflecting an angle across the $x$-axis leaves the horizontal coordinate
unchanged. Since cosine measures the horizontal coordinate of a point on
the unit circle, $\cos(-\theta) = \cos\theta$.
\end{proof}

These two properties are traditionally called odd symmetry and even
symmetry, and each is a direct geometric consequence of what sine and
cosine measure on the circle, rather than an independent fact requiring
separate proof.

\section{The Tangent Function}

Recall from Chapter~4 that $\tan\theta = \sin\theta/\cos\theta$, undefined
whenever $\cos\theta = 0$. This occurs at $\pi/2$, $3\pi/2$, $5\pi/2$, and
every angle coterminal with them, so the graph of tangent contains a
vertical asymptote at each of these values; unlike sine and cosine, the
tangent graph is unbounded, stretching from $-\infty$ to $\infty$ between
each pair of consecutive asymptotes.

\begin{theorem}
The tangent function has period $\pi$.
\end{theorem}

\begin{proof}
Adding $\pi$ to an angle reverses the sign of both sine and cosine:
$\sin(\theta+\pi) = -\sin\theta$ and $\cos(\theta+\pi) = -\cos\theta$.
Therefore
\[
\tan(\theta+\pi) = \frac{-\sin\theta}{-\cos\theta} = \frac{\sin\theta}{\cos\theta} = \tan\theta,
\]
so the graph of tangent repeats every $\pi$ units, half the period of sine
and cosine, precisely because the two sign changes that occur after adding
only $\pi$ cancel inside the ratio.
\end{proof}

\section{Transformations of Trigonometric Graphs}

Most applications involve a modified trigonometric function rather than the
basic sine or cosine curve, typically written
\[
y = A\sin(B\theta - C) + D.
\]
Each parameter controls a distinct geometric feature of the graph. The
factor $A$ scales the graph vertically, so that the amplitude becomes
$|A|$. The factor $B$ compresses or stretches the graph horizontally,
changing the period from $2\pi$ to $2\pi/|B|$, since a full cycle of the
inner angle $B\theta - C$ now requires $\theta$ to advance through only
$2\pi/|B|$ units rather than a full $2\pi$. The quantity $C/B$ shifts the
graph horizontally and is called the \emph{phase shift}, since setting
$B\theta - C = 0$ shows that the graph's usual starting point has been
delayed until $\theta = C/B$. Finally, the constant $D$ moves the entire
graph vertically, so that its new midline is $y = D$ rather than $y = 0$.

\begin{algorithmproc}{How to Graph a Trigonometric Function from Its Parameters}
Given $y = A\sin(B\theta - C) + D$, first determine the amplitude $|A|$ and
the period $2\pi/|B|$. Next determine the phase shift $C/B$, taking care
with its sign, and the vertical shift $D$, which fixes the new midline.
Use these four values to sketch one complete cycle of the curve using its
key points, then extend the sketch periodically in both directions.
\end{algorithmproc}

\begin{example}
Graph $y = 3\sin\left(2\theta - \dfrac{\pi}{2}\right) + 1$. The amplitude is
$3$, the period is $2\pi/2 = \pi$, the phase shift is $(\pi/2)/2 = \pi/4$
to the right, and the vertical shift is $1$, giving a midline of $y=1$. The
graph therefore oscillates between $4$ and $-2$, completing one full cycle
every $\pi$ units, with that cycle beginning $\pi/4$ units later than an
unshifted sine curve would.
\end{example}

\begin{practicenow}
Determine the amplitude, period, phase shift, and vertical shift of
$y = 2\cos(3\theta + \pi) - 4$.
\end{practicenow}

\section{Applications}

The graphs developed in this chapter model an enormous range of periodic
phenomena, from seasonal temperature variation and ocean tides to sound
waves, alternating electrical current, and biological rhythms. The same
curves that emerge from a single point rotating around the unit circle in
Chapter~5 turn out to describe some of the most important oscillatory
processes in the natural and engineered world, which is exactly the payoff
promised at the end of that chapter.

\section{Looking Ahead}

The graphs of sine and cosine reveal an important complication: a
horizontal line typically intersects either graph at infinitely many
points, so that an equation such as $\sin\theta = 1/2$ has infinitely many
solutions rather than one. Sine and cosine, as they stand, therefore cannot
be reversed uniquely across their entire domains, and if a trigonometric
function is to be "undone" at all, its domain must first be restricted to a
piece on which it is one-to-one. The next chapter takes up this problem
directly, developing the inverse trigonometric functions and showing how an
angle can be recovered from a known ratio once the appropriate restriction
has been made.

\section{Exercises}
\begin{exercise}
State the amplitude, period, phase shift, and vertical shift of
$y = 4\sin\left(\dfrac{\theta}{2} - \pi\right) + 2$, and sketch one cycle.
\end{exercise}

\begin{exercise}
Explain, using the proof of the tangent period theorem, why the period of
$\tan\theta$ is $\pi$ while the period of $\sin\theta$ and $\cos\theta$ is
$2\pi$.
\end{exercise}

\begin{exercise}
Find all solutions of $\cos\theta = -\dfrac{\sqrt2}{2}$ in the interval
$[0, 4\pi)$, and use the graph of cosine to explain why there is more than
one solution.
\end{exercise}

\begin{exercise}
Write an equation of the form $y = A\sin(B\theta - C) + D$ for a curve with
amplitude $5$, period $4\pi$, no phase shift, and midline $y = -3$.
\end{exercise}
